%! TEX root = ../edtc-18-19.tex

\chapter{Ecuații diferențiale de ordinul întîi}

\indent\indent {\color{red} Dacă nu se precizează altfel, vom presupune că lucrăm 
  cu funcții de forma $y = y(x)$, deci $y'$ va însemna $\dfrac{dy}{dx}$.}

\section{Ecuații cu variabile separabile/separate}
\label{sec:sep}
\index{ecuații!cu variabile separabile}

\indent\indent Acesta este cel mai simplu exemplu de ecuații diferențiale și se 
rezolvă direct prin integrare, după o reordonare corespunzătoare.

\textbf{Exemplu 1}: $(1 + x^2) yy' + x(1 + y^2) = 0$, știind că $y(1) = 2$.

\emph{Soluție:} Separăm variabilele și diferențialele și obținem succesiv:
\begin{align*}
  (1 + x^2) y \cdot \frac{dy}{dx} + x(1 + y^2) &= 0 \Leftrightarrow\\
  (1 + x^2) y dy &= -x(1 + y^2) dx \Leftrightarrow \\
  \frac{y}{1 + y^2} dy &= -\frac{x}{1 + x^2} dx \Leftrightarrow \\
  \frac{1}{2} \ln (y^2 + 1) &= -\frac{1}{2} \ln(x^2 + 1) + c.
\end{align*}

Pentru uniformitate, putem pune $\dfrac{1}{2} \ln c$ în locul constantei.

Rezultă $1 + y^2 = \dfrac{c}{1 + x^2}$, cu $c > 0$, pentru existența logaritmului.

Înlocuim în condiția inițială $y(1) = 2$ și obținem $c = 10$. În fine:
\[
  y(x) = \sqrt{\frac{9 - x^2}{1 + x^2}}, \quad x \in (-3, 3),
\]
punînd și condițiile de existență.

\vspace{2cm}

În unele cazuri, este util să facem schimbări de variabilă. De exemplu:

\textbf{Exemplu 2}: $y' = \sin^2(x - y)$.

\emph{Soluție}: Notăm $x - y = z$ și avem că $y' = 1 - z'$, de unde, în
ecuație, obținem succesiv:
\begin{align*}
  z' & = 1 - \sin^2 z = \cos^2 z \Rightarrow\\
  \frac{dz}{dx} &= \cos^2 z \Rightarrow \\
  \frac{dz}{\cos^2 z} &= dx \Rightarrow \\
  \tan z &= x + c \Rightarrow \\
  \tan (x - y) &= x + c,
\end{align*}
care poate fi prelucrată pentru a obține $y(x)$ sau lăsată în forma implicită.

%%%%%%%%%%%%%%%%%%%%%%%%%%%%%%%%%%%%%%%%%%%%%%%%%%%%%%%%%%%%%%%%%%%%%%

\section{Ecuații liniare}
\label{sec:lin}
\index{ecuații!liniare}

Forma generală a acestor ecuații este:
\[
  y' = P(x) \cdot y + Q(x).
\]
Distingem două cazuri:
\begin{itemize}
\item Dacă $Q = 0$, atunci ecuația se numește \emph{omogenă};
\item Dacă $Q \neq 0$, ecuația este \emph{neomogenă}.
\end{itemize}

Metoda generală de rezolvare este să folosim formula:
\[
  y(x) = c \cdot \exp \Big( -\int_{x_0}^x P(t) dt \Big),
\]
pentru a obține o \emph{soluție particulară pentru ecuația omogenă}. Apoi, din teorie,
știm că putem folosi \emph{metoda variației constantelor (Lagrange)} pentru a obține
soluția ecuației neomogene. Pentru aceasta, în locul constantei $c$ vom considera
o funcție $c(x)$ și înlocuim în ecuația inițială.

Pe scurt, pentru ecuația liniară:
\[
  y' + P(x) y = Q(x),
\]
avem:
\begin{itemize}
  \item Soluția particulară pentru varianta omogenă $y_p = c \cdot \exp\Big( -\int P(x) dx \Big)$;
  \item Soluția generală (din Lagrange):
    \[
      y_g = \exp\Big(-\int P(x) dx\Big) \cdot \Big( C + \int Q(x) \cdot \exp\Big( \int P(x) dx \Big) dx \Big)
    \]
\end{itemize}

Soluția generală a problemei este suma între soluția particulară și cea generală.


\textbf{Exemplu 1}: $y' + y \sin x = -\sin x \cos x$.

\emph{Soluție}: Putem înlocui direct în formulă și obținem $y = C \cdot e^{\cos x} - \cos x - 1$.

\vspace{1cm}

\textbf{Exemplu 2}: $y' + xy = x e^{-\frac{x^2}{2}}$.

\emph{Soluție}: Asociem ecuația omogenă $y' + xy = 0$, pe care o rescriem:
\[
  \frac{dy}{dx} + xy = 0 \Rightarrow \frac{dy}{y} = -xdx \Rightarrow y_p = ce^{-\frac{x^2}{2}}.
\]

Acum, folosind metoda lui Lagrange, căutăm o soluție de forma $y(x) = c(x) e^{-\frac{x^2}{2}}$.

Înlocuim în ecuația inițială și găsim $y' + xy = c'(x) e^{-\frac{x^2}{2}}$, dar, comparînd
cu ecuația dată, găsim $c'(x) = x$, de unde $c(x) = \dfrac{x^2}{2}$.

Așadar, avem $y_g = \dfrac{x^2}{2} e^{-\frac{x^2}{2}}$, iar soluția finală este suma celor două:
\[
  y = \Big( \frac{x^2}{2} + c \Big) e^{-\frac{x^2}{2}}.
\]

\vspace{1cm}

\textbf{Exemplu 3}: $2xyy' + 2y^2 - x^4 = 0$.

\emph{Soluție:} Putem face o substituție $y^2 = z$, cu care
ecuația devine $z' + \frac{2}{x} z = x^3$. Avem, în acest caz, $P(x) = \frac{2}{x}$,
iar $Q(x) = x^3$. Cum $\int P(x) dx = 2 \ln x$, iar $\int Q(x) e^{\int P(x) dx} dx = \frac{x^6}{6}$,
obținem soluția generală:
\[
  y(x) = \frac{1}{x} \sqrt{c + \frac{x^6}{6}}, x > 0.
\]

%%%%%%%%%%%%%%%%%%%%%%%%%%%%%%%%%%%%%%%%%%%%%%%%%%%%%%%%%%%%%%%%%%%%%%

\section{Ecuația Bernoulli}
\label{sec:bernoulli}
\index{ecuații!Bernoulli}

Forma generală a ecuației Bernoulli este:
\[
  y' + P(x) y = Q(x) y^\alpha.
\]

Remarcăm că:
\begin{itemize}
\item Dacă $\alpha = 0$, obținem o ecuație omogenă;
\item Dacă $\alpha \neq 0$, obținem o ecuație neomogenă, ca în secțiunea anterioară.
\end{itemize}

Pașii de rezolvare sînt:
\begin{itemize}
\item Se împarte cu $y^\alpha$ și obținem:
  \[
    y^{-\alpha} y' + P(x) y^{1 - \alpha} = Q(x);
  \]
\item Facem substituția $y^{1 - \alpha} = z$ și ajungem la ecuația:
  \[
    (1 - \alpha) y^{1 - \alpha} \cdot y' = z',
\]
de unde $\dfrac{z'}{1 - \alpha} + P(x) z = Q(x)$, care este o ecuație
neomogenă, rezolvabilă ca în secțiunea anterioară.
\end{itemize}

\textbf{Exemplu 1}: $y' - \dfrac{y}{3x} = \dfrac{1}{3} y^4 \ln x, x > 0$.

\emph{Soluție:} Avem $\alpha = 4$, deci împărțim la $y^4$ și obținem:
\[
  y^{-4} y' - \frac{1}{3x} y^{-3} = \frac{1}{3} \ln x.
\]

Cu substituția $z = y^{-3}$, ajungem la:
\[
  z' = - 3y^4 y' \Rightarrow z' + \frac{1}{x} z = -\ln x.
\]
Avem $P(x) = \frac{1}{x}$ și $Q(x) = -\ln x$, deci putem aplica formula pentru
soluția generală a ecuației neomogene:
\[
  z = e^{-\ln x} \Big( c - \int \ln x e^{\ln x} dx \Big) = \frac{1}{x} \Big( c - \int x \ln x \Big).
\]

În fine:
\[
  y^{-3} = \frac{c}{x} + \frac{x}{4} - \frac{x}{2} \ln x, x > 0.
\]

\vspace{1cm}

\textbf{Exemplu 2}: $y' + \dfrac{2}{3x}y = \dfrac{1}{3} y^2$.

\emph{Soluție}: Avem $\alpha = 2$, deci împărțim la $y^2$ și ajungem la:
\[
  \frac{y'}{y^2} + \frac{2}{3x} \cdot \frac{1}{y} = \frac{1}{3}.
\]

Cu substituția $z = \dfrac{1}{y}$, avem $z' + \dfrac{2}{3x} z = \dfrac{1}{3}$,
care este liniară și neomogenă.

Lucrăm cu ecuația omogenă $z' + \dfrac{2}{3x} = 0$, de unde $\dfrac{z'}{z} = \dfrac{2}{3x}$,
care este cu variabile separabile și găsim $z = cx^{\frac{2}{3}}$, \emph{pentru ecuația omogenă}.

Aplicăm acum metoda variației constantelor și luăm $z = c(x) x^{\frac{2}{3}}$. După înlocuire
în ecuația inițială, avem $c(x) = -x^{\frac{1}{3}}$.

Soluția finală este acum suma $z = -x + cx^{\frac{2}{3}} = \dfrac{1}{y}$.

%%%%%%%%%%%%%%%%%%%%%%%%%%%%%%%%%%%%%%%%%%%%%%%%%%%%%%%%%%%%%%%%%%%%%%

\section{Ecuația Riccati}
\label{sec:ec-ric}
\index{ecuații!Riccati}

Forma generală este:
\[
  y' = P(x) y^2 + Q(x) y + R(x).
\]

Observăm că:
\begin{itemize}
\item Dacă $Q = 0$, avem o ecuație liniară și neomogenă;
\item Dacă $R = 0$, este o ecuație Bernoulli, cu $\alpha = 2$.
\end{itemize}
  
În general, ecuația Riccati se rezolvă știind o soluție particulară $y_p$. Apoi,
folosind formula $y = y_p + z$ și înlocuind, se ajunge la o ecuație Bernoulli:
\[
  z' + \big(p(x) + 2q(x) y_p\big)z + q(x) z^2 = 0.
\]

Pentru a găsi soluții particulare, următoarea teoremă ne ajută într-un anumit caz:

\begin{theorem}\label{thm:ric}
  Dacă avem ecuația Riccati de forma:
  \[
    y' = Ay^2 + \frac{B}{x}y + \frac{C}{x^2},
  \]
  unde $(B + 1)^2 - 4AC \geq 0$, atunci o soluție particulară este $y_p = \dfrac{1}{x}$.
\end{theorem}

\textbf{Exemplu 1}: $y' = -\dfrac{1}{3} y^2 - \dfrac{2}{3x^2}, x > 0$.

\emph{Soluție}: Aplicînd teorema, putem lua $y_p = \dfrac{1}{x}$ ca soluție particulară.
Apoi căutăm o soluție generală de forma $y = \dfrac{1}{x} + z$, dar, pentru conveniență,
putem lua $z \to \dfrac{1}{z}$. Înlocuim în ecuație și obținem succesiv:
\begin{align*}
  -\frac{1}{x^2} - \frac{z'}{z^2} &= -\frac{1}{3} \Big( \frac{1}{x^2} + \frac{2}{xz} + \frac{1}{z^2} \Big) - \frac{2}{3x} \Rightarrow \\
  z' - \frac{2}{3x} z &= \frac{1}{3} \Rightarrow \\
  z &= Cx^{\frac{2}{3}} + x,
\end{align*}
de unde $y = \dfrac{1}{x} + \dfrac{1}{Cx^{\frac{2}{3}} + x}$, cu $x > 0$.

%%%%%%%%%%%%%%%%%%%%%%%%%%%%%%%%%%%%%%%%%%%%%%%%%%%%%%%%%%%%%%%%%%%%%%

\section{Ecuația Clairaut}
\label{sec:clairaut}
\index{ecuații!Clairaut}

Forma generală a ecuației este:
\[
  y = xy' + \varphi(y').
\]

Pentru rezolvare, se notează $y' = p$ și, înlocuind, avem:
\[
  y = xp + \varphi(p).
\]

Derivăm după $x$ și găsim:
\[
  p = p + x \cdot \frac{dp}{dx} + \varphi'(p) \cdot \frac{dp}{dx},
\]
de unde $\big( x + \varphi'(p) \big) \frac{dp}{dx} = 0$.

Mai departe:
\begin{itemize}
\item Dacă $\dfrac{dp}{dx} = 0$, atunci $p = C$ și avem $y = C \cdot x + \varphi(c)$, ca soluție generală;
\item Dacă $x + \varphi'(p) = 0$, obținem \emph{soluția singulară}, care se prezintă parametric, adică
  în funcție de $p$, astfel:
  \[
    \begin{cases}
      x &= -\varphi'(p) \\
      y &= -p\varphi'(p) + \varphi(p)
    \end{cases}
  \]
\end{itemize}

\vspace{1cm}

\textbf{Exemplu:} $y = xy' - {y'}^2$.

\emph{Soluție:} Notăm $y' = p$, deci $y = xp - p^2$.

Obținem, înlocuind în ecuație:
\[
  pdx + xdp - 2pdp = pdx \Rightarrow (x - 2p) dp = 0.
\]

Distingem cazurile:
\begin{itemize}
\item Dacă $dp = 0$, atunci $p = C$ și $y = Cx - C_2$, o soluție particulară;
\item Soluția singulară se reprezintă parametric, corespunzător cazului $x - 2p = 0$, prin:
  \[
    \begin{cases}
      x &= 2p \\
      y &= p^2
    \end{cases}
  \]
\end{itemize}

Revenind la $y$, găsim $y = \dfrac{x^2}{4}$.

\begin{remark}\label{rk:infas}
  Geometric, soluția generală este \emph{înfășurătoarea} soluției particulare, adică o curbă
  care aproximează, din aproape în aproape, dreapta soluției particulare.
\end{remark}

%%%%%%%%%%%%%%%%%%%%%%%%%%%%%%%%%%%%%%%%%%%%%%%%%%%%%%%%%%%%%%%%%%%%%%

\section{Ecuații exacte. Factor integrant}
\label{sec:exc-int}
\index{ecuații!exacte}
\index{ecuații!exacte!factor integrant}

Forma generală este:
\[
  P(x, y) + Q(x, y) y' = 0.
\]

Dacă are loc $\dfrac{\partial P}{\partial y} = \dfrac{\partial Q}{\partial x}$, atunci
ecuația se numește \emph{exactă}, iar soluția ei generală se prezintă prin:
\[
  F(x, y) = \int_{x_0}^x P(t, y_0) dt + \int_{y_0}^y Q(x, t) dt,
\]
pentru orice $(x, y)$ din domeniul de definiție, iar $(x_0, y_0)$ un punct arbitrar fixat,
în jurul căruia determinăm soluția.

\begin{remark}\label{rk:partial}
  Pentru simplitate, vom mai folosi notațiile $P_y = \dfrac{\partial P}{\partial y}$ și
  similar pentru $Q_x$. Deci condiția de exactitate se mai scrie, pe scurt, $P_y = Q_x$.
\end{remark}

\vspace{1cm}

\textbf{Exemplu:} $(3x^2 - y) + (3y^2 - x) y' = 0$.

\emph{Soluție:} Remarcăm că avem $P(x, y) = 3x^2 - x$ și $Q(x, y) = 3y^2 - x$, deci
$P_y = Q_x = -1$, ecuația fiind exactă. Obținem direct, din formula pentru soluția generală:
\[
  F(x, y) = x^3 + y^3 - xy + x_0y_0 - x_0^3 - y_0^3.
\]

Cum $(x_0, y_0)$ sînt fixate, putem prezenta soluția în \emph{forma implicită} $y = \varphi$,
unde $\varphi = x^3 + y^3 - xy = c$, constanta fiind expresia de $(x_0, y_0)$ de mai sus.

\vspace{1cm}

Dacă ecuația nu este exactă, se caută \emph{un factor integrant}, adică o funcție $\mu(x, y)$
cu care se înmulțește ecuația pentru a deveni exactă.

Există două cazuri particulare în care factorul integrant se găsește simplu:
\begin{itemize}
\item Dacă expresia $\dfrac{P_y - Q_x}{Q}$ depinde doar de $x$, se poate lua $\mu = \mu(x)$;
\item Dacă expresia $\dfrac{Q_x - P_y}{P}$ depinde doar de $y$, se poate lua $\mu = \mu(y)$.
\end{itemize}

În celelalte cazuri, factorul integrant, de obicei, se dă în enunțul problemei, deoarece este
dificil de găsit, dar o teoremă de existență ne arată că el poate fi mereu găsit.

\begin{remark}\label{rk:int}
  Pentru a reține mai ușor condițiile simple de căutare a factorului integrant, pornim cu
  ecuația în forma generală, înmulțim cu $\mu$, ceea ce ne schimbă și $P$ și $Q$, apoi scriem
  condiția de exactitate. Pentru primul caz, de exemplu, obținem:
  \[
    \frac{\mu_x}{\mu} = \frac{P_y - Q_x}{Q},
  \]
  deci dacă membrul drept depinde doar de $x$, putem lua $\mu = \mu(x)$ și similar în al doilea caz.
\end{remark}

Odată găsit factorul integrant, să presupunem $\mu = \mu(x)$, integrăm condiția de exactitate și găsim:
\[
  \ln \mu(x) = \int \varphi(x) + c,
\]
unde $\varphi(x) = \dfrac{P_y - Q_x}{Q}$ și deci $\mu(x) = \exp \Big( \int \varphi(x) dx \Big)$.

\vspace{1cm}

\textbf{Exemplu}: $(1 - x^2 y) + x^2 (y - x) y' = 0$.

\emph{Soluție:} Observăm că ecuația nu este exactă, dar verificînd condițiile pentru factorul
integrant, putem lua $\mu(x) = \dfrac{1}{x^2}$. Înmulțim ecuația cu $\mu$ și găsim:
\[
  \Big( \frac{1}{x^2} - y \Big) + (y - x)y' = 0,
\]
care este exactă. Atunci putem scrie direct soluția generală:
\[
  F(x, y) = \frac{y^2}{2} - \frac{1}{x} - xy + C,
\]
soluția implicită fiind $\dfrac{y^2}{2} - \dfrac{1}{x} - xy = \text{const.}$.

\vspace{1cm}

\textbf{Exercițiu:} Rezolvați ecuația $y^2 (2x - 3y) + (7 - 3xy^2) y' = 0$, căutînd
(după verificare) un factor integrant $\mu(y)$.

%%%%%%%%%%%%%%%%%%%%%%%%%%%%%%%%%%%%%%%%%%%%%%%%%%%%%%%%%%%%%%%%%%%%%%

\section{Ecuația Lagrange}
\label{sec:lagr}
\index{ecuații!Lagrange}

Forma generală a ecuației Lagrange este:
\[
  A(y')\cdot y + B(y')\cdot x + C(y') = 0.
\]

Dacă avem $A(y') \neq 0$, putem împărți și obținem:
\[
  y = f(y') + g(y'), \quad f(y') = \frac{- B(y')}{A(y')}, g(y') = -\frac{C(y')}{A(y')}.
\]

Dacă $f(y') = y'$, obținem o ecuație de tip Clairaut.

Altfel, fie $y' = p$, deci $dy = pdx$. Înlocuim în ecuația inițială și găsim:
\[
  y = xf(p) + g(p) \Rightarrow dy = f(p) dx + xf'(p) dp + g'(p) dp.
\]

Egalăm cele două expresii pentru $dy$ și ajungem la:
\[
  pdx = f(p) dx + xf'(p) dp + g'(p) dp.
\]

În fine:
\[
  \big( f(p) - p \big) dx + \big( xf'(p) + g'(p) \big) dp = 0.
\]

În acest caz, dacă $f(p)$ este o constantă, obținem o ecuație cu variabile separabile.

Altfel, putem împărți la $(f(p) - p) dp$ și ajungem la:
\[
  \frac{dx}{dp} + \frac{f'(p)}{f(p) - p} x + \frac{g'(p)}{f(p) - p} = 0.
\]

Aceasta este o ecuație liniară și neomogenă în $x(p)$ și o rezolvăm corespunzător.

Dacă $f(p) = p$, obținem o soluție particulară.

\vspace{1cm}

\textbf{Exemplu:} $y = x - \dfrac{4}{9} {y'}^2 + \dfrac{8}{27} {y'}^3$.

\emph{Soluție:} Facem notația $y' = p$, deci $dy = pdx$. Obținem:
\[
  y = x - \frac{4}{9} p^2 + \frac{8}{27} p^3 \Rightarrow %
  dy = dx - \frac{8}{9} pdp + \frac{8}{9} p^2 dp.
\]

Rezultă $pdx = dx - \dfrac{8}{9} pdp + \dfrac{8}{9} p^2 dp$, deci:
\[
  (p - 1)(dx - \frac{8}{9} pdp) = 0.
\]

Distingem cazurile:

Dacă $dx = \dfrac{8}{9} pdp$, atunci $x = \dfrac{4}{9} p^2 + c$ și, înlocuind în
ecuația inițială, găsim $y = \dfrac{8}{27} p^3 + c$. Soluția poate fi prezentată parametric:
\[
  \begin{cases}
    x &= \dfrac{4}{9} p^2 + c \\
      & \\
    y &= \dfrac{8}{27} p^3 + c
  \end{cases}
\]

Dacă $p = 1$, avem o soluție particulară $y = x + c$, iar constanta se obține a fi $c = -\dfrac{4}{27}$,
din ecuația inițială.


%%%%%%%%%%%%%%%%%%%%%%%%%%%%%%%%%%%%%%%%%%%%%%%%%%%%%%%%%%%%%%%%%%%%%%

\section{Exerciții}

Rezolvați următoarele ecuații diferențiale, precizînd și tipul lor:
\begin{enumerate}[(1)]
  \item $xy^\prime - y + 2x^2 y^2 = 0 $;
  \item $ (3x^2 + 2)y^\prime + xy^2 = 0 $;
  \item $(2x - y + 2) dx + (-x + y + 1) dy = 0 $;
  \item $ xy^\prime = 4y + x^2 \sqrt{y} $;
  \item $ (x^2 + 1) y^\prime + 2xy^2 = 0 $;
  \item $ (2x + y + 1) dx + (x + y - 1) dy = 0 $;
  \item $ y^\prime = \dfrac{x^2}{y} $;
  \item $ y^\prime = \dfrac{y}{x} + x $;
  \item $ 3x^2 y dx - (\ln y + x^3) dy = 0 $;
  \item $ y^\prime = x^2 y $;
  \item $ y^\prime = \dfrac{x^2 + y^2}{2xy} $;
  \item $ y^\prime = 2xy^2 $;
  \item $ y^\prime = \dfrac{4x}{y^2} $;
  \item $ (x + y + 1)dx + (x - y^3 + 3) dy = 0 $;
  \item $ y^\prime = \dfrac{1}{x}y - \dfrac{\ln x}{x} $;
  \item $ y^\prime = -\dfrac{y}{x} + \dfrac{e^x}{x} $;
  \item $ (x + y^2) dx + y(1 - x)dy = 0 $;
  \item $ 3x^2 y dx - (x^3 + \ln y) dy = 0 $;
  \item $ y^\prime = -2y + x^2 + 2x $;
  \item $ y^\prime = \dfrac{2y}{x} + \dfrac{3}{x^2} $;
  \item $ y^\prime = 2xy + x^3 y^{\dfrac{1}{2}} $;
  \item $ y^\prime = -\dfrac{1}{x}y + \dfrac{\ln x}{x} y^2 $.
\end{enumerate}

